% GENERATED FILE. Edit canonical lessons, then run npm run generate:books.
% canonical-source-hash: fnv1a64:9b7b1d4d598e71de
% canonical-lessons: MW-C133-chana, MW-C133-mug, MW-C133-moth, MW-C133-til, MW-C133-gur

\chapter{Chickpeas, Mung beans, Moth beans, Sesame, Jaggery}
\label{ch:mw-chickpeas-mung-beans-moth-beans-sesame-jaggery}
\hlchaptermodality{133}

\begin{chapteropening}
\textbf{By the end of this chapter:} \emph{I can say chickpeas, mung beans, moth beans, sesame and jaggery.}

Complete the last lesson of chapter 133: I can say chickpeas, mung beans, moth beans, sesame and jaggery.
\end{chapteropening}

\section[\texorpdfstring{\mw{चणा} (chaṇā)}{चणा (chaṇā)}]{\mw{चणा} (chaṇā) — chickpeas}

\label{lesson:MW-C133-chana}



\begin{quote}
Before the new one: say the Marwadi for maize, then the Marwadi for wheat.
\end{quote}

\subsection*{You'll want to know: \mw{चणा}}
\textbf{\mw{चणा}} — \emph{chaṇā} — \textquotedblleft{}chickpeas\textquotedblright{}.

\begin{cousinweb}
\textbf{\mw{चणा}} (\emph{chaṇā}) is a tadbhava, a word inherited from Sanskrit and reshaped by sound change. It comes from \emph{caṇa}, \textquotedblleft{}the chickpea\textquotedblright{}. CDIAL gives no Marwari form in this entry (CDIAL 4579); the match is through Gujarati \emph{caṇā}. Marwadi, like Gujarati, keeps the \emph{ṇ}.
\end{cousinweb}

\begin{grammarlens}[title={What you've built}]
One more word for this chapter.
\end{grammarlens}

\subsection*{Your turn}
\begin{itemize}
\raggedright
  \item \emph{Say it:} \emph{chaṇā}
  \item \emph{Say it:} \emph{chaṇā}, once more
  \item \emph{Say it:} say \emph{chaṇā}
  \item \emph{Recall:} say the Marwadi for maize, then the Marwadi for wheat, then say \emph{chaṇā} again
\end{itemize}

\begin{tcolorbox}[breakable,colback=teal!4,colframe=teal!35!black,title={Before you move on}]
What does \mw{चणा} mean? (\textquotedblleft{}Chickpeas\textquotedblright{}.) Say it once more.
\end{tcolorbox}
\section[\texorpdfstring{\mw{मूंग} (mū̃g)}{मूंग (mū̃g)}]{\mw{मूंग} (mū̃g) — mung beans}

\label{lesson:MW-C133-mug}



\begin{quote}
Before the new one: say the Marwadi for wheat, then the Marwadi for chickpeas.
\end{quote}

\subsection*{You'll want to know: \mw{मूंग}}
\textbf{\mw{मूंग}} — \emph{mū̃g} — \textquotedblleft{}mung beans\textquotedblright{}.

\begin{cousinweb}
\textbf{\mw{मूंग}} (\emph{mū̃g}) is a tadbhava, a word inherited from Sanskrit and reshaped by sound change. It comes from \emph{mudga}, \textquotedblleft{}the bean Phaseolus mungo\textquotedblright{}. CDIAL gives no Marwari form in this entry (CDIAL 10198); the match is through Hindi \emph{mū̃g}.
\end{cousinweb}

\begin{grammarlens}[title={What you've built}]
One more word for this chapter.
\end{grammarlens}

\subsection*{Your turn}
\begin{itemize}
\raggedright
  \item \emph{Say it:} \emph{mū̃g}
  \item \emph{Say it:} \emph{mū̃g}, once more
  \item \emph{Say it:} say \emph{mū̃g}
  \item \emph{Recall:} say the Marwadi for wheat, then the Marwadi for chickpeas, then say \emph{mū̃g} again
\end{itemize}

\begin{tcolorbox}[breakable,colback=teal!4,colframe=teal!35!black,title={Before you move on}]
What does \mw{मूंग} mean? (\textquotedblleft{}Mung beans\textquotedblright{}.) Say it once more.
\end{tcolorbox}
\section[\texorpdfstring{\mw{मोठ} (moṭh)}{मोठ (moṭh)}]{\mw{मोठ} (moṭh) — moth beans}

\label{lesson:MW-C133-moth}



\begin{quote}
Before the new one: say the Marwadi for chickpeas, then the Marwadi for mung beans.
\end{quote}

\subsection*{You'll want to know: \mw{मोठ}}
\textbf{\mw{मोठ}} — \emph{moṭh} — \textquotedblleft{}moth beans\textquotedblright{}.

\begin{cousinweb}
\textbf{\mw{मोठ}} (\emph{moṭh}) is a tadbhava, a word inherited from Sanskrit and reshaped by sound change. It comes from \emph{mukuṣṭha}, \textquotedblleft{}a kind of bean\textquotedblright{}. CDIAL gives no Marwari form in this entry (CDIAL 10148); the match is through Hindi \emph{moṭh}. The Sanskrit word has many variant spellings in the old word-lists.
\end{cousinweb}

\begin{grammarlens}[title={What you've built}]
One more word for this chapter.
\end{grammarlens}

\subsection*{Your turn}
\begin{itemize}
\raggedright
  \item \emph{Say it:} \emph{moṭh}
  \item \emph{Say it:} \emph{moṭh}, once more
  \item \emph{Say it:} say \emph{moṭh}
  \item \emph{Recall:} say the Marwadi for chickpeas, then the Marwadi for mung beans, then say \emph{moṭh} again
\end{itemize}

\begin{tcolorbox}[breakable,colback=teal!4,colframe=teal!35!black,title={Before you move on}]
What does \mw{मोठ} mean? (\textquotedblleft{}Moth beans\textquotedblright{}.) Say it once more.
\end{tcolorbox}
\section[\texorpdfstring{\mw{तिल} (til)}{तिल (til)}]{\mw{तिल} (til) — sesame}

\label{lesson:MW-C133-til}



\begin{quote}
Before the new one: say the Marwadi for mung beans, then the Marwadi for moth beans.
\end{quote}

\subsection*{You'll want to know: \mw{तिल}}
\textbf{\mw{तिल}} — \emph{til} — \textquotedblleft{}sesame\textquotedblright{}.

\begin{cousinweb}
\textbf{\mw{तिल}} (\emph{til}) is a tadbhava, a word inherited from Sanskrit and reshaped by sound change. It comes from \emph{tila}, \textquotedblleft{}sesame\textquotedblright{}. CDIAL lists Old Marwari \emph{tila} under it (CDIAL 5827). The Old Marwari row means \textquotedblleft{}small particle\textquotedblright{}, a sense the Sanskrit word also had. CDIAL thinks the Sanskrit word itself probably came from Munda.
\end{cousinweb}

\begin{grammarlens}[title={What you've built}]
One more word for this chapter.
\end{grammarlens}

\subsection*{Your turn}
\begin{itemize}
\raggedright
  \item \emph{Say it:} \emph{til}
  \item \emph{Say it:} \emph{til}, once more
  \item \emph{Say it:} say \emph{til}
  \item \emph{Recall:} say the Marwadi for mung beans, then the Marwadi for moth beans, then say \emph{til} again
\end{itemize}

\begin{tcolorbox}[breakable,colback=teal!4,colframe=teal!35!black,title={Before you move on}]
What does \mw{तिल} mean? (\textquotedblleft{}Sesame\textquotedblright{}.) Say it once more.
\end{tcolorbox}
\section[\texorpdfstring{\mw{गुड़} (guṛ)}{गुड़ (guṛ)}]{\mw{गुड़} (guṛ) — jaggery}

\label{lesson:MW-C133-gur}



\begin{quote}
Before the new one: say the Marwadi for moth beans, then the Marwadi for sesame.
\end{quote}

\subsection*{You'll want to know: \mw{गुड़}}
\textbf{\mw{गुड़}} — \emph{guṛ} — \textquotedblleft{}jaggery\textquotedblright{}.

\begin{cousinweb}
\textbf{\mw{गुड़}} (\emph{guṛ}) is a tadbhava, a word inherited from Sanskrit and reshaped by sound change. It comes from \emph{guḍa}, \textquotedblleft{}boiled sugarcane juice, molasses\textquotedblright{}. CDIAL gives no Marwari form in this entry (CDIAL 4182); the match is through Hindi \emph{guṛ}. CDIAL asks whether it is the same word as \emph{guḍa}, \textquotedblleft{}ball\textquotedblright{}.
\end{cousinweb}

\begin{grammarlens}[title={What you've built}]
One more word for this chapter.
\end{grammarlens}

\subsection*{Your turn}
\begin{itemize}
\raggedright
  \item \emph{Say it:} \emph{guṛ}
  \item \emph{Say it:} \emph{guṛ}, once more
  \item \emph{Say it:} say \emph{guṛ}
  \item \emph{Recall:} say the Marwadi for moth beans, then the Marwadi for sesame, then say \emph{guṛ} again
\end{itemize}

\begin{tcolorbox}[breakable,colback=teal!4,colframe=teal!35!black,title={Before you move on}]
What does \mw{गुड़} mean? (\textquotedblleft{}Jaggery\textquotedblright{}.) Say it once more.
\end{tcolorbox}
